\label{chap:appendix}

\section{Mathematical Proofs of Nilpotent Triangular Norm Properties}

\subsection{Proof of Triangular Norm Axioms for \L{}ukasiewicz \texorpdfstring{$t$}{t}-Norm}
\begin{Theorem}
The function $T_L(a, b) = \max(0, a + b - 1)$ is a valid continuous triangular norm on $[0, 1] \times [0, 1] \to [0, 1]$.
\end{Theorem}
\begin{proof}
We verify the four axioms:
\begin{enumerate}
    \item \textbf{Commutativity:} $T_L(a, b) = \max(0, a + b - 1) = \max(0, b + a - 1) = T_L(b, a)$ for all $a, b \in [0, 1]$.
    \item \textbf{Boundary Conditions:} For any $a \in [0, 1]$:
    \begin{align}
    T_L(a, 1) &= \max(0, a + 1 - 1) = \max(0, a) = a \quad (\text{since } a \ge 0) \\
    T_L(a, 0) &= \max(0, a + 0 - 1) = \max(0, a - 1) = 0 \quad (\text{since } a \le 1 \implies a - 1 \le 0)
    \end{align}
    \item \textbf{Monotonicity:} If $a \le c$ and $b \le d$, then $a + b - 1 \le c + d - 1$. Since $f(z) = \max(0, z)$ is non-decreasing on $\mathbb{R}$:
    \begin{equation}
    T_L(a, b) = \max(0, a + b - 1) \le \max(0, c + d - 1) = T_L(c, d)
    \end{equation}
    \item \textbf{Associativity:} For all $a, b, c \in [0, 1]$:
    \begin{equation}
    T_L(T_L(a, b), c) = \max\left(0, \max(0, a + b - 1) + c - 1\right) = \max(0, a + b + c - 2)
    \end{equation}
    By symmetry, $T_L(a, T_L(b, c)) = \max(0, a + b + c - 2)$, establishing associativity.
\end{enumerate}
All four axioms hold; hence, $T_L$ is a valid continuous $t$-norm.
\end{proof}

\subsection{Unit Gradients on Active Regions}
For any latent attribute coordinate $d$ on the active domain where $Q_{id} + K_{jd} - 1 - \tau > 0$:
\begin{equation}
\mathcal{M}_{ijd} = Q_{id} + K_{jd} - 1 - \tau \implies \frac{\partial \mathcal{M}_{ijd}}{\partial Q_{id}} = 1.0, \quad \frac{\partial \mathcal{M}_{ijd}}{\partial K_{jd}} = 1.0
\end{equation}
Whenever the coordinate sum exceeds $1 + \tau$, the gradient with respect to either input is constant: $\frac{\partial}{\partial Q} = 1.0$ and $\frac{\partial}{\partial K} = 1.0$. This means active connections do not suffer from the vanishing or exploding gradients often caused by softmax exponents.

\section{Automated Unit Testing Suite (19 Passing Tests)}
We wrote 19 automated unit tests using `pytest` to check that our implementation maintains core architectural invariants—including bounded embeddings, operator boundary conditions, and KV-cache numerical equivalence (Table \ref{tab:unit_test_suite}). All 19 tests pass.

\begin{table}[h]
\centering
\caption{Automated Unit Testing Suite (19 Passing Tests)}
\label{tab:unit_test_suite}
\small
\resizebox{\columnwidth}{!}{%
\begin{tabular}{lp{6.8cm}c}
\toprule
\textbf{Test Module} & \textbf{Property Verified} & \textbf{Status} \\
\midrule
\texttt{test\_embeddings.py} & Representations bounded in $[0, 1]^D$, convex blending closure & \textbf{PASSED} (5/5) \\
\texttt{test\_lattice\_ops.py} & $t$-norm axioms, nilpotent zero boundaries, unit gradients & \textbf{PASSED} (5/5) \\
\texttt{test\_kv\_cache.py} & KV-cache numerical equivalence ($\Delta = 5.06 \times 10^{-7}$), diagonal fallback & \textbf{PASSED} (3/3) \\
\texttt{test\_encoder.py} & Bidirectional symmetric routing, MLM head, distractor isolation & \textbf{PASSED} (3/3) \\
\texttt{test\_model\_forward.py} & Exact parameter parity against baseline, loss shape verification & \textbf{PASSED} (3/3) \\
\midrule
\textbf{Total Test Suite} & \textbf{19 Independent Unit Tests} & \textbf{19 / 19 PASSED} \\
\bottomrule
\end{tabular}%
}
\end{table}

\section{Reproducibility Protocol and Environment Setup}
To reproduce our results, clone the repository and run the following commands with `uv`:
\begin{verbatim}
uv sync                                                 # Sync environment
uv run pytest tests/                                    # Run 19 unit tests
uv run python generate.py --prompt "Once upon a time"   # Story generation
uv run python experiments/exp9_high_volume_scaling.py   # 6.14M token scaling
uv run python experiments/exp10_encoder_edge.py         # Encoder MLM test
\end{verbatim}
